\begin{beginnerstatement}[Kurz erklärt: Besondere Projekttypen]

Native Integration bedeutet, besondere Funktionen eines Betriebssystems oder
Geräts direkt zu verwenden.

Bei Hard Real Time muss eine Reaktion unter allen vorgesehenen Bedingungen
eine feste Zeitgrenze einhalten.

Eine Game Engine stellt spezielle Werkzeuge für grafikintensive Spiele bereit.

Ein Embedded System ist ein Rechner in einem Gerät; seine hardwarenahe
Software heißt Firmware.

Eine native Mobile App wird gezielt für ein mobiles Betriebssystem entwickelt,
während HPC besonders rechenintensive Aufgaben auf leistungsstarken Systemen
bezeichnet.

Compliance bedeutet, verbindliche rechtliche, vertragliche oder
organisatorische Vorgaben einzuhalten und dies bei Bedarf nachzuweisen.

Für Projekte mit einem solchen technischen Kern liefert das Template keine
vollständige Basis; bei hohen Compliance-Anforderungen sind zusätzliche
Prüfungen nötig.

\end{beginnerstatement}
